\chapter[2020 --- Alter et al.]{2020: Harvey J. Alter, Michael Houghton, Charles M. Rice}
\label{ch:2020_Alter}

\section*{Main Contribution}

\textit{Discovered Hepatitis C virus, enabling blood testing and antiviral treatments that have saved millions of lives and brought the world close to eliminating the disease.}\cite{nobel-2020,lecture-2020}

\section{Official Citation}

for their discovery of Hepatitis C virus

\section{Background and Historical Context}

The 2020 Nobel Prize in Physiology or Medicine recognized Harvey J. Alter, Michael Houghton, and Charles M. Rice ``for their discovery of Hepatitis C virus.'' Their work solved one of the major puzzles in modern medicine---the identification of the previously unknown virus responsible for the majority of post-transfusion hepatitis cases---and opened the door to the development of diagnostic tests, antiviral therapies, and ultimately the possibility of eradicating hepatitis C virus (HCV) infection worldwide.

Viral hepatitis, a group of inflammatory diseases of the liver caused by hepatitis viruses, has been one of a significant global health challenges of the twentieth and twenty-first centuries. Five major hepatitis viruses have been identified: hepatitis A virus (HAV), hepatitis B virus (HBV), hepatitis C virus (HCV), hepatitis D virus (HDV), and hepatitis E virus (HEV). Of these, HBV and HCV are the most dangerous, causing chronic infections that can lead to cirrhosis, liver failure, and hepatocellular carcinoma (liver cancer).

The identification of HBV by Baruch Blumberg in the 1960s was a major milestone in the fight against viral hepatitis. Blumberg discovered the ``Australia antigen'' (later identified as the hepatitis B surface antigen) and developed a blood test that could screen blood donations for HBV infection. The introduction of HBV screening in the early 1970s dramatically reduced the incidence of post-transfusion hepatitis B.

However, the elimination of post-transfusion hepatitis B did not eliminate the problem. A large proportion of post-transfusion hepatitis cases continued to occur even after the implementation of HBV screening, indicating the existence of a second agent---later designated ``non-A, non-B hepatitis'' (NANBH). This mysterious agent was responsible for the majority of post-transfusion hepatitis cases and was a major cause of chronic liver disease, cirrhosis, and liver cancer.

The identification of the NANBH agent was one of the major and challenging problems in virology and hepatology during the 1970s and 1980s. The difficulty in identifying the agent was compounded by its apparent inability to be cultured in cell culture or to be transmitted to laboratory animals in a reliable and reproducible manner. The development of new molecular approaches was needed to identify and characterize this elusive pathogen.

\section{The Discovery/Contribution}

The discovery of HCV was achieved through the complementary efforts of three investigators: Harvey Alter, who provided the clinical and epidemiological framework; Michael Houghton, who identified the virus using molecular cloning; and Charles Rice, who proved that the cloned virus could cause disease.

\textbf{Harvey J. Alter:} Alter, an American physician-scientist born in 1935 in New York City, working at the National Institutes of Health (NIH), was a pioneer in the study of post-transfusion hepatitis. Beginning in the 1960s, Alter conducted a series of epidemiological studies that established the existence and clinical significance of NANBH. He developed a chimpanzee model of NANBH that enabled the study of the disease and provided a system for testing potential diagnostic and therapeutic approaches.

Alter's epidemiological studies demonstrated that NANBH was a distinct clinical entity that was different from hepatitis A and hepatitis B. He showed that NANBH was transmitted through blood and blood products, that it frequently led to chronic infection, and that chronic NANBH was a significant cause of cirrhosis and liver cancer. These observations established the public health importance of NANBH and created a sense of urgency about identifying the responsible agent.

Alter also made the critical observation that NANBH could be transmitted to chimpanzees, providing an animal model for studying the disease. This model was essential for the subsequent identification and characterization of HCV, as it enabled the testing of the cloned virus for its ability to cause disease.

\textbf{Michael Houghton:} Houghton, a British virologist born in 1956 in the United Kingdom, working at Chiron Corporation (later acquired by Novartis) in Emeryville, California, used molecular cloning to identify the previously unknown agent responsible for NANBH. Houghton's approach was based on the hypothesis that if the NANBH agent was a virus, it should contain nucleic acid (either DNA or RNA) that could be detected using molecular hybridization techniques.

Houghton and his colleagues developed a novel and ambitious strategy for identifying the NANBH virus. They started with a chimpanzee that was chronically infected with NANBH, extracted RNA from its liver tissue, and constructed a complementary DNA (cDNA) library---a collection of DNA clones representing all of the RNA sequences present in the infected liver. They then screened this library using serum from patients with NANBH, looking for cDNA clones that reacted with antibodies in the patient serum. This approach was based on the assumption that if a viral sequence was present in the cDNA library, the immune system of infected patients would have produced antibodies against the viral protein encoded by that sequence.

After extensive screening of thousands of cDNA clones, Houghton and his colleagues identified a single clone that reacted specifically with serum from NANBH patients but not with serum from uninfected individuals. This clone, which they designated 5-1-1, was shown to encode part of the genome of a previously unknown RNA virus, which they named hepatitis C virus. The discovery was reported in a landmark paper published in \textit{Science} in 1989.

The identification of HCV was a breakthrough of enormous significance. For the first time, the mysterious agent responsible for the majority of post-transfusion hepatitis had been identified and molecularly characterized. The availability of the viral genome enabled the development of diagnostic blood tests that could screen blood donations for HCV infection, dramatically reducing the risk of transfusion-transmitted hepatitis C. It also enabled the development of molecular tools for studying the virus and its pathogenesis.

\textbf{Charles M. Rice:} Rice, an American virologist born in 1952 in Sacramento, California, working at Washington University in St. Louis, made the critical contribution of proving that the cloned HCV could cause disease. While the identification of HCV by Houghton was a landmark achievement, some scientists remained skeptical that the cloned virus was truly the cause of NANBH, as Koch's postulates had not been fully satisfied. The cloned virus had been identified based on its immunoreactivity with patient antibodies, but it had not been shown to be capable of causing disease when introduced into a susceptible host.

Rice addressed this challenge by constructing a full-length, infectious clone of the HCV genome and demonstrating that it could cause hepatitis when inoculated into chimpanzees. Rice and his colleagues synthesized a complete HCV genome from overlapping cDNA clones and showed that RNA transcribed from this clone was infectious---it could initiate a productive viral infection in chimpanzees that led to the development of hepatitis with biochemical and histological features characteristic of NANBH.

Rice's experiment provided the definitive proof that HCV was the causative agent of NANBH, fulfilling Koch's postulates and establishing the causal relationship between the virus and the disease. This achievement was essential for the acceptance of HCV as a distinct hepatitis virus and for the development of antiviral therapies targeting the virus.

Rice also made important contributions to understanding the molecular biology of HCV, including the characterization of its polyprotein processing, the identification of its non-structural proteins, and the development of cell culture systems for studying the viral life cycle. These advances enabled the development of antiviral drugs targeting specific steps in the HCV replication cycle.

\section{Impact on Modern Medicine}

The discovery of HCV has had a transformative impact on medicine, leading to the development of diagnostic tests, antiviral therapies, and prevention strategies that have saved millions of lives.

\textbf{Diagnostic tests:} The identification of the HCV genome enabled the development of blood screening tests that could detect HCV antibodies and viral RNA in blood donations. The introduction of HCV screening in the early 1990s dramatically reduced the risk of transfusion-transmitted hepatitis C, from approximately 10\% of transfused patients to near zero. Screening tests have also enabled the identification of individuals with chronic HCV infection, many of whom are asymptomatic and unaware of their infection, allowing early diagnosis and treatment.

\textbf{Antiviral therapy:} The understanding of the HCV molecular biology, built upon the foundation laid by Alter, Houghton, and Rice, has enabled the development of direct-acting antiviral (DAA) drugs that can cure chronic HCV infection with cure rates exceeding 95\%. The development of DAAs was a notable achievements in the history of antiviral therapy. Early treatments for hepatitis C relied on pegylated interferon-alpha and ribavirin, which had modest efficacy (cure rates of approximately 40--50\%) and significant side effects. The development of DAAs, beginning with the approval of sofosbuvir (Sovaldi) in 2013, revolutionized the treatment of hepatitis C.

Sofosbuvir, a nucleotide analog inhibitor of the HCV NS5B RNA-dependent RNA polymerase, was a breakthrough in hepatitis C treatment. In clinical trials, sofosbuvir-based regimens achieved cure rates exceeding 90\% with treatment durations of 8--12 weeks and minimal side effects. The development of additional DAAs targeting other HCV proteins, including NS3/4A protease inhibitors (simeprevir, paritaprevir, glecaprevir), NS5A inhibitors (ledipasvir, daclatasvir, velpatasvir, elbasvir), and NS5B non-nucleoside inhibitors (dasabuvir), has provided a range of combination regimens that can cure virtually all genotypes of HCV.

\textbf{Public health impact:} The availability of highly effective DAA therapy has transformed hepatitis C from a chronic, potentially fatal disease into a curable condition. The World Health Organization (WHO) has set the goal of eliminating hepatitis C as a public health threat by 2030, and many countries have implemented national elimination strategies based on universal screening, treatment with DAAs, and prevention measures.

However, the high cost of DAAs has limited access to treatment in many low- and middle-income countries, where the burden of hepatitis C is notable. Efforts to reduce the cost of DAAs and to expand access to treatment through generic drug production and voluntary licensing agreements are ongoing and are critical for achieving the WHO elimination goal.

Harvey J. Alter, Michael Houghton, and Charles M. Rice were awarded the Nobel Prize in Physiology or Medicine in 2020 ``for their discovery of Hepatitis C virus.'' Their work has fundamentally changed the understanding and treatment of hepatitis C and has opened the door to the eventual elimination of this devastating disease.

\section{Legacy}

The legacy of Alter, Houghton, and Rice's work extends across multiple areas of medicine and public health. Their discovery of HCV solved one of the major puzzles in modern medicine and provided the foundation for the development of diagnostic tests, antiviral therapies, and prevention strategies that have saved millions of lives.

Harvey Alter's epidemiological studies established the clinical significance of NANBH and provided the framework for the subsequent identification of HCV. His work demonstrated the importance of careful clinical observation and epidemiological investigation in the identification of new infectious agents.

Michael Houghton's use of molecular cloning to identify HCV was a pioneering application of molecular biology to the discovery of new pathogens. His approach has been adapted for the identification of other previously unknown viruses and has become a standard methodology in the field of pathogen discovery.

Charles Rice's demonstration that the cloned HCV could cause disease was a landmark achievement in virology. His work fulfilled Koch's postulates and established the causal relationship between HCV and hepatitis, providing the definitive proof needed for the acceptance and clinical development of HCV-targeted therapies.

Together, the work of these three laureates has transformed the landscape of viral hepatitis research and has brought the world closer to the goal of eliminating hepatitis C as a public health threat. Their legacy continues to inspire new research and clinical innovation, and their contributions will remain central to the field of hepatology for generations to come.


\section{Modern Developments}

Alter, Houghton, and Rice's HCV work has led to modern hepatology. Direct-acting antivirals (sofosbuvir, ledipasvir) cure >95% of HCV infections in 8-12 weeks. Understanding of HCV has enabled blood screening, preventing transfusion-transmitted infection. Understanding of HCV vaccine development remains challenging but active. Understanding of liver fibrosis has led to research on anti-fibrotic therapies for cirrhosis. The goal of HCV elimination by 2030 is achievable with current tools.


\section{Bibliography}

\begin{thebibliography}{9}
\bibitem{nobel-2020}
Nobel Prize Outreach, ``The Nobel Prize in Physiology or Medicine 2020,''\emph{NobelPrize.org}.\\
\url{https://www.nobelprize.org/prizes/medicine/2020/summary/}

\bibitem{lecture-2020}
Harvey J. Alter, ``Nobel Lecture,''\emph{NobelPrize.org}. Winner-authored primary source; downloadable from the official Nobel Prize archive.\\
\url{https://www.nobelprize.org/prizes/medicine/2020/alter/lecture/}
\end{thebibliography}
